\section{练习}

\begin{question}
\begin{drill}[单透镜的像位置与放大率]\label{d:31}
空气中薄凸透镜焦距 \(10.0\,\mathrm{cm}\)，高 \(2.00\,\mathrm{mm}\) 的直立小物在镜前 \(30.0\,\mathrm{cm}\)。求像的位置、高度、实虚和正倒，并用两条辅助光线核对。
\end{drill}
\fig{fig:d31q}{练习：单透镜的像位置与放大率}{
\begin{tikzpicture}
\draw[axisline](-3.5,0)--(2.5,0);\draw[<->,>=Stealth,world](0,-.9)--(0,.9);
\draw[axis,thick](-3,0)--(-3,.5)node[above]{$B$};
\node[below]at(-3,0){$A$};\node[below]at(0,0){$O$};
\fill(1,0)circle(1pt)node[below]{$F'$};\fill(-1,0)circle(1pt)node[below]{$F$};
\end{tikzpicture}}
\end{question}


\begin{question}
\begin{drill}[两薄透镜依次成像]\label{d:32}
空气中两薄透镜同轴，\(f_1=+10.0\,\mathrm{cm}\)、\(f_2=-10.0\,\mathrm{cm}\)，间隔 \(15.0\,\mathrm{cm}\)。实物在第一镜前 \(20.0\,\mathrm{cm}\)。求中间像、第二镜的物距、最终像及总横向放大率。
\end{drill}
\fig{fig:d32q}{练习：两薄透镜依次成像}{
\begin{tikzpicture}
\draw[axisline](-2.6,0)--(3.2,0);
\draw[<->,>=Stealth,world](0,-.8)--(0,.8)node[above]{$L_1$};
\draw[>-<,>=Stealth,world](1.5,-.8)--(1.5,.8)node[above]{$L_2$};
\draw[axis,thick](-2,0)--(-2,.4)node[above]{$B$};
\node[below]at(-2,0){$A$};\draw[dim](0,-1.1)--(1.5,-1.1)node[midway,below]{$15.0\,\mathrm{cm}$};\draw[dim](-2,-1.8)--(0,-1.8)node[midway,below]{$20.0\,\mathrm{cm}$};
\end{tikzpicture}}
\end{question}


\begin{question}
\begin{drill}[光焦度分配与消色差]\label{d:33}
两紧贴薄透镜的总焦距要求为 \(10.0\,\mathrm{cm}\)。从红光到蓝光，两材料透镜的光焦度变化分别满足 \(\Delta\Phi_C=0.0200\Phi_C\)、\(\Delta\Phi_F=0.0400\Phi_F\)，其中光焦度带符号。求满足一阶消色差的两片焦距。
\end{drill}
\fig{fig:d33q}{练习：光焦度分配与消色差}{
\begin{tikzpicture}
\draw[axisline](-2,0)--(3,0);
\draw[world](0,-1)--(0,1)node[above]{$C$};
\draw[world](.6,-1)--(.6,1)node[above]{$F$};
\node[below,align=center]at(.3,-1.4){紧贴组合\\总光焦度 \(0.100\,\mathrm{cm^{-1}}\)};
\end{tikzpicture}}
\end{question}

